\chapter{Optical Product Readiness: From Requirements to Fleet}
\label{ch:product-readiness}
\label{ch:validation}

Read this chapter for: phase-label translation, evidence-discipline
boundaries, the 11-step product-readiness lifecycle, essential bring-up and
system-corner logic, and readiness interview practice.

Use the detailed chapters for: link physics, source and wavelength mechanisms,
reliability qualification, manufacturing validation, networking, and failure
analysis.

Use the measurement-reference appendix for: instruments, metric procedures,
stressed-receiver methods, detailed CMIS guidance, and link-budget accounting
conventions
(\Cref{app:measurement-reference}).

An optical product does not become ready through one activity called validation.
It becomes ready through several distinct evidence disciplines.
Characterization maps behavior. Verification checks frozen requirements.
System validation demonstrates intended use. Reliability qualification
addresses permanent degradation from time and exposure. Manufacturing
validation demonstrates that the production system can reproduce and control
the design. Production acceptance testing screens units or populations using
validated measurements and proxies. Pilot and fleet evidence compare the
release model with operation.

This chapter explains how an optical product progresses from requirements and
architectural assumptions to verified performance, validated system use,
qualified reliability, controlled manufacturing, deployment, and fleet
learning
(\Cref{sec:product-readiness,app:decision-trees,sec:tree-evidence-block}).

\keyidea{Product readiness and validation are not synonyms.
Product readiness is the complete evidence-building process used to move an
optical product from requirements and architectural assumptions through
characterization, requirement verification, system validation, reliability
qualification, manufacturing validation, controlled deployment, and fleet
learning. System validation is one discipline within that lifecycle. It
demonstrates that the complete product is suitable for its intended use with
the supported hosts, peers, fiber plant, management behavior, environmental
conditions, and operational workload. This chapter therefore uses
\emph{product-readiness lifecycle} for the complete 11-step flow and reserves
\emph{validation} for the narrower engineering claim.}

Companies may organize these activities under phase names such as EVT, DVT, or
PVT, but those labels do not replace the technical definitions. Organizational
team names are not used as engineering definitions in this book. Abbreviations
are collected in \Cref{app:abbreviations}.

\section{Why program-phase names are not enough}
\label{sec:phase-labels}

Hardware programs often use EVT\sidenote{EVT: Engineering Validation Test or
Engineering Verification Test, depending on the organization. EVT is a
program-phase label, not a universal engineering definition. It commonly
emphasizes early hardware, architecture learning, bring-up, characterization,
and retirement of high-risk assumptions. Ask which hardware revision,
evidence, and exit decision the program expects.},
DVT\sidenote{DVT: commonly Design Validation Test, but some organizations
use Design Verification Test. The label may include requirement verification,
system validation, interoperability, margin testing, and reliability evidence
in different proportions. This book uses the precise terms verification and
system validation for the engineering claims rather than relying on DVT
alone.},
PVT\sidenote{PVT: commonly Production Validation Test or Production
Verification Test. It usually emphasizes production-intent hardware, factory
measurements, ATP coverage, process capability, traceability, supplier
readiness, and controlled ramp. A PVT label does not by itself establish
system suitability or lifetime qualification.},
qualification, pilot, and production-readiness labels as though they were
universal technical definitions. They are not. They are company-specific
containers for work. One organization may perform reliability qualification
during DVT, while another begins it in EVT and completes it during PVT. One
may call a build ``DVT'' when it is primarily verifying requirements; another
may use the same label for interoperability and system validation.

The stable engineering questions are more useful than the phase names: What
behavior has been characterized? Which requirements have been verified? Has
intended system use been validated? Have lifetime mechanisms been qualified?
Can the production system reproduce and control the design?

\keyidea{When someone says EVT, DVT, or PVT, ask four questions.
Which hardware and software configuration is included? Which engineering
evidence is expected? Which risks remain intentionally open? What decision or
exit gate does the phase enable? The phase name is useful only after those
four questions are answered.}

After pilot, programs often speak of
MP\sidenote{MP: mass production. MP describes production exposure and
operating cadence, not evidence quality by itself. A product is ready for
unrestricted MP only when the applicable design, reliability, manufacturing,
supply, service, and monitoring decisions have been supported.}
(mass production) and sustaining. Those labels still need an exit decision.

\begin{table*}[ht]
\refstepcounter{table}\label{tab:phase-readiness-map}%
\centering
\setlength{\tabcolsep}{3pt}%
\footnotesize
\renewcommand{\arraystretch}{1.15}%
\begin{tabular}{@{}p{0.20\linewidth}p{0.28\linewidth}p{0.46\linewidth}@{}}
\toprule
Program label & Common emphasis & Typical product-readiness work \\
\midrule
EVT & Architecture and engineering learning & Architecture review, prototype
bring-up, risk retirement, early characterization, measurement development \\
\addlinespace[0.6em]
DVT & Design evidence and intended-use closure & Characterization completion,
requirement verification, margin, interoperability, system validation,
substantial qualification evidence \\
\addlinespace[0.6em]
PVT & Production-intent and ramp readiness & Production-reference freeze,
manufacturing validation, measurement-system analysis, ATP coverage, process
capability, controlled ramp evidence \\
\addlinespace[0.6em]
Pilot / limited deployment & Bounded operational evidence & Cohort deployment,
enhanced telemetry, success criteria, rollback, fleet comparison \\
\addlinespace[0.6em]
MP / sustaining & Controlled volume and field learning & Ramp, SPC, supplier
controls, fleet monitoring, failure analysis, next-revision feedback \\
\bottomrule
\end{tabular}
\par\medskip
{\raggedright\footnotesize\textbf{Table~\thetable.} Illustrative mapping from
program-phase labels to product-readiness work. EVT, DVT, and PVT are not
standards, and their activities overlap. Always ask what hardware population,
evidence, and exit decision the phase label represents. Reliability planning
may begin in EVT, representative qualification may run through DVT, and
corrective requalification may continue into PVT.\par}
\end{table*}

\section{The evidence disciplines and their decisions}
\label{sec:evidence-disciplines}

Organizational ownership varies. The technical question should remain stable
even when one engineer or team performs several disciplines.

\emph{Characterization} maps nominal behavior, distributions, sensitivities,
interactions, and failure boundaries. It asks what the design does rather than
only whether it passes.

\emph{Verification} compares measured or analyzed evidence with a frozen
requirement under stated conditions and at a named reference plane. Ask: does
the evidence meet the specified requirement?

\emph{System validation} demonstrates that the complete product is suitable for
its intended system use, including supported hosts, peers, fiber plants,
firmware, management behavior, environmental conditions, and operational
interactions. Ask: does the complete product work for the intended use and
supported ecosystem?

\emph{Reliability qualification} builds bounded confidence that credible
lifetime and environmental mechanisms will not cause unacceptable permanent
degradation during the intended use period. Ask: will the product continue
meeting its requirements after the intended time and exposure?

\emph{Manufacturing validation} demonstrates that the production system can
repeatedly build, measure, trace, detect, and control the design supported by
the earlier engineering evidence. Ask: can the factory reproduce and control
the design at the required quality and scale?

\emph{Production acceptance testing}, commonly called
ATP\sidenote{ATP: acceptance test procedure or acceptance test plan,
depending on company usage. This book uses ATP for the released production
acceptance-test flow that makes a pass, fail, hold, rework, or disposition
decision using validated measurements or proxies. ATP does not replace
characterization, system validation, reliability qualification, or
manufacturing validation.}, applies validated production measurements or
proxies to make a disposition decision for an individual unit or production
population. Ask: should this unit or population proceed?

\emph{Pilot deployment}\sidenote{Pilot: a bounded production-representative
deployment with an identifiable cohort, explicit success criteria, enhanced
telemetry, and a rollback or containment plan. A pilot is not merely a small
shipment.} asks whether a bounded production-representative population behaves
as predicted operationally.

\emph{Fleet monitoring} compares deployed behavior with the release assumptions
and identifies changes by lane, module, lot, site, firmware, topology, and
installation age. Ask: is the deployed population behaving as predicted?

\begin{table*}[ht]
\refstepcounter{table}\label{tab:validation-jobs}%
\centering
\setlength{\tabcolsep}{2.5pt}%
\footnotesize
\renewcommand{\arraystretch}{1.12}%
\begin{tabular}{@{}p{0.18\linewidth}p{0.26\linewidth}p{0.26\linewidth}p{0.24\linewidth}@{}}
\toprule
Discipline & Primary question & Main output & What it does not replace \\
\midrule
Characterization & How does the design behave and vary? & Behavioral model,
distributions, sensitivities, failure boundaries & Requirement verification or
release approval \\
\addlinespace[0.5em]
Verification & Does measured evidence meet the frozen requirement? & Traceable
requirement result with conditions and uncertainty & Intended-use evidence \\
\addlinespace[0.5em]
System validation & Is the complete product suitable for its intended system
use? & Supported operating and interoperability envelope & Lifetime
qualification or factory capability \\
\addlinespace[0.5em]
Reliability qualification & Will time and exposure cause unacceptable permanent
degradation? & Mechanism-specific life and environmental confidence & Every-unit
screening or manufacturing reproducibility \\
\addlinespace[0.5em]
Manufacturing validation & Can the production system repeatedly build, measure,
trace, and control the design? & Production capability, measurement confidence,
control plan, ramp evidence & System validation or life qualification \\
\addlinespace[0.5em]
Production acceptance testing & Should this unit or population proceed? & Pass,
fail, hold, rework, or disposition record & Complete mechanism coverage or
design qualification \\
\addlinespace[0.5em]
Pilot deployment & Does a bounded production-representative population behave as
predicted operationally? & Controlled field evidence and expansion decision &
Uncontrolled mass deployment \\
\addlinespace[0.5em]
Fleet monitoring & Does the deployed population match the release model? &
Trends, cohorts, alerts, incident evidence, next-revision learning & Root-cause
confirmation by itself \\
\bottomrule
\end{tabular}
\par\medskip
{\raggedright\footnotesize\textbf{Table~\thetable.} Evidence disciplines and
the decisions they can honestly support. Abbreviations:
\Cref{app:abbreviations}.\par}
\end{table*}

The same measurement can contribute to more than one discipline. What changes
is the question, population, condition, acceptance criterion, and decision. A
temperature sweep may characterize reversible performance, verify a
maximum-temperature requirement, support system validation in a loaded chassis,
or provide pre- and post-stress measurements for qualification. The equipment
does not determine the engineering category; the claim does.

The product-readiness lifecycle in \Cref{sec:product-readiness} sequences these
disciplines so expensive evidence is not asked to answer the wrong question.

\section{The optical product-readiness lifecycle}
\label{sec:product-readiness}
\label{sec:validation-workflow}

\refstepcounter{table}\label{tab:ladder}\label{tab:ladder-gates}%
\endgraf\medskip\noindent
\fcolorbox{keyidearule}{codebg}{%
\begin{minipage}{\dimexpr\textwidth-2\fboxsep-2\fboxrule\relax}
{\raggedright\footnotesize\textbf{Table~\thetable.} Optical product-readiness
lifecycle at a glance (Steps 1--11). Wall-chart:
\Cref{sec:tree-qualification}. Interview drill:
\Cref{sec:interview-validation-ladder}. NPI names: \Cref{tab:npi}.\par}
\medskip
\textbf{The product-readiness lifecycle removes a different uncertainty at each
step. The steps are logically ordered, although much of the work overlaps in
calendar time.}
{\raggedright\footnotesize
\begin{itemize}
\item \mbox{\textbf{Step 1: Define requirements.}} Uncertainty: what must the
      product do, and under what conditions? Decision: establish the evidence
      contract.
\item \mbox{\textbf{Step 2: Review architecture.}} Uncertainty: can the proposed
      design plausibly close the budgets and risks? Decision: proceed, modify,
      or reject the architecture.
\item \mbox{\textbf{Step 3: Bring up hardware.}} Uncertainty: is there a stable,
      known, reproducible operating state? Decision: begin interpretable
      characterization.
\item \mbox{\textbf{Step 4: Characterize behavior.}} Uncertainty: how does the
      design behave and vary? Decision: select corners and freeze evidence
      targets.
\item \mbox{\textbf{Step 5: Verify and validate.}} Uncertainty: does it meet
      requirements and work in the intended system? Decision: define the
      supported operating envelope. This step alone does not make the product
      ready for unrestricted production.
\item \mbox{\textbf{Step 6: Qualify reliability.}} Uncertainty: will time and
      exposure cause unacceptable permanent degradation? Decision: support the
      life and environmental claim.
\item \mbox{\textbf{Step 7: Validate manufacturing.}} Uncertainty: can
      production reproduce and control the design? Decision: authorize
      controlled production exposure.
\item \mbox{\textbf{Step 8: Run controlled pilot.}} Uncertainty: does a bounded
      deployed population behave as predicted? Decision: expand, hold,
      restrict, or roll back.
\item \mbox{\textbf{Step 9: Ramp mass production.}} Uncertainty: do factory,
      supplier, and field indicators remain controlled as exposure grows?
      Decision: increase production volume.
\item \mbox{\textbf{Step 10: Monitor fleet.}} Uncertainty: does deployed
      behavior match the release model? Decision: contain, sustain, or
      investigate.
\item \mbox{\textbf{Step 11: Feed learning forward.}} Uncertainty: which
      assumptions, controls, or architectures must change? Decision: improve
      the next revision.
\end{itemize}
\textit{Each step answers a question that the previous step cannot. A
later-stage success cannot compensate for missing evidence from an earlier
stage.}\par}
\end{minipage}%
}%
\endgraf\medskip

Program names such as
NPI\sidenote{NPI: New Product Introduction. NPI is the cross-functional program
that moves a design into controlled production. It includes engineering
evidence, supplier readiness, factory process development, test development,
quality planning, operations, change control, and ramp. NPI is an
organizational umbrella, not a technical evidence category.}
organize the same work without replacing the step definitions
(\Cref{tab:npi,sec:phase-labels}).

\subsection{Logical order does not require calendar serialization}

The steps are ordered by the questions they answer, not by a rule that one
department must finish before another begins. Architecture review, test
development, supplier qualification, reliability planning, factory preparation,
and telemetry design often overlap in calendar time. The important requirement
is that later decisions do not claim evidence that earlier work has not
established. Parallel execution is healthy. Ambiguous evidence ownership is not.

\subsection{How the lifecycle changes}

Before hardware, requirements and architecture define success and decide whether
the budgets can close under stated assumptions. With hardware, bring-up makes
measurements interpretable; characterization maps behavior; verification and
system validation close frozen requirements and intended use. After the shipping
envelope is understood, reliability qualification, manufacturing validation,
pilot, ramp, fleet monitoring, and feedback remove the remaining uncertainties
that no quiet-bench close can answer.

\begin{table*}[ht]
\refstepcounter{table}\label{tab:readiness-ownership}%
\centering
\setlength{\tabcolsep}{3pt}%
\footnotesize
\begin{tabular}{@{}p{0.48\linewidth}p{0.46\linewidth}@{}}
\toprule
Topic & Primary owner \\
\midrule
IM/DD architecture and link budget & \Cref{ch:imdd} \\
Noise, sensitivity, RIN, BER, and metric interpretation & \Cref{ch:models} \\
Sources and modulation & \Cref{ch:lasers} \\
WDM and wavelength control & \Cref{ch:wdm} \\
Product-readiness sequence and system-validation decisions &
\Cref{ch:product-readiness} \\
Reliability qualification & \Cref{ch:reliability} \\
Manufacturing validation and production control & \Cref{ch:manufacturing} \\
Networking, FEC, retrains, telemetry, and availability & \Cref{ch:networking} \\
Failure analysis and recurrence control & \Cref{ch:failure-modes} \\
Instruments, procedures, and test-reference material &
\Cref{app:measurement-reference} \\
Decision navigation & \Cref{app:decision-trees} \\
\bottomrule
\end{tabular}
\par\medskip
{\raggedright\footnotesize\textbf{Table~\thetable.} Where detailed ownership
lives. Chapter~7 keeps the readiness sequence and decisions.\par}
\end{table*}

\section{Steps 1--11}
\label{sec:readiness-steps}

\subsection{Step 1: Define the requirements}
\label{sec:ladder-requirements}

Define what success means before instruments enter the conversation. Write
performance, environment, reliability, manufacturing, and operational
requirement classes with owners and named planes where applicable
(\Cref{sec:systems-engineering-loop,tab:laser-prd}).

\textit{Evidence and handoff:} Freeze a signed requirements slice specific
enough for architecture and later pass criteria; refuse hardware spend until
success is defined.

\subsection{Step 2: Review the architecture}
\label{sec:ladder-architecture}

Decide whether the architecture can meet the requirements before tooling makes
changes expensive. Close optical, noise, thermal, electrical, reliability, and
manufacturing budgets on stated assumptions
(\Cref{sec:systems-engineering-loop,tab:laser-engineering-checklist}).

\textit{Evidence and handoff:} Proceed to bring-up only when budgets close or
open items are named with redesign triggers; otherwise redesign first.

\subsection{Step 3: Bring up the hardware}
\label{sec:ladder-bringup}

Establish a known, reproducible operating state so later sweeps are
interpretable. Separate integration fails from product-performance questions.
Detail lives in \Cref{sec:bringup}.

\textit{Evidence and handoff:} Continue to characterization when identity,
rails, management-ready state, basic optical power, lock, and a simple link are
reproducible; otherwise debug integration.

\subsection{Step 4: Characterize the behavior}
\label{sec:ladder-characterization}

Build the behavioral model: nominals, distributions, sensitivities, and cliffs.
A characterization cliff improves understanding; it does not automatically fail
the product (\Cref{sec:evidence-disciplines}).

\textit{Evidence and handoff:} Name thin ledgers and candidate corners; proceed
to Step~5, derate, or redesign before loaded-fleet claims.

\subsection{Step 5: Verify requirements and validate system use}
\label{sec:ladder-margin}
\label{sec:ladder-margin-interop}
\label{sec:ladder-interop}

Verification closes frozen requirements at named planes with stated uncertainty.
System validation closes intended use across supported hosts, peers, firmware,
fiber plants, and workloads. Related evidence can serve both claims; the claim
still differs (\Cref{sec:evidence-disciplines}).

\paragraph{Verify requirements and margin.}
\label{sec:ladder-verify}
Margin testing is verification when it checks a defined margin requirement under
specified conditions. Track optical, electrical, thermal, and control ledgers
separately and stack once. Do not mix average-power and OMA budgets, reference
planes, or embedded and explicit transmitter penalties
(\Cref{sec:link-budget}).

\paragraph{Validate interoperability and intended use.}
\label{sec:ladder-system-validation}
Exercise production-representative hosts, peers, firmware, chassis loading, and
plant conditions. Golden-host margin is not the ecosystem exit
(\Cref{sec:prod-corners}).

Receiver-margin evidence may include a named stressed-receiver method,
including SECQ where applicable to the PMD; the method, stress calibration, and
metric must be stated explicitly (\Cref{sec:secq}).

Transmitter evidence should include average power and modulation-quality
evidence such as OMA, RLM, and the applicable transmitter-quality metric.
Passing average power does not establish a valid PAM4 transmitter. A composite
transmitter metric supports acceptance or margin accounting but does not
identify a unique physical mechanism (\Cref{sec:tdecq,ch:models,ch:lasers}).

\begin{table*}[ht]
\refstepcounter{table}\label{tab:evidence-domains}%
\centering
\setlength{\tabcolsep}{3pt}%
\footnotesize
\begin{tabular}{@{}p{0.14\linewidth}p{0.26\linewidth}p{0.28\linewidth}p{0.26\linewidth}@{}}
\toprule
Evidence domain & Readiness question & Representative evidence & Detailed
ownership \\
\midrule
Transmitter & Is sufficient modulated optical quality launched? & OMA, level
quality, wavelength, applicable Tx-quality metric &
\Cref{ch:models,ch:lasers,ch:wdm,app:measurement-reference} \\
\addlinespace[0.5em]
Channel & Does the supported plant preserve required margin? & Insertion loss,
ORL/reflection, dispersion or filtering &
\Cref{ch:imdd,ch:models,ch:wdm} \\
\addlinespace[0.5em]
Receiver & Can the receiver meet the stated objective under named stress? &
Sensitivity, overload, stressed-receiver evidence &
\Cref{ch:imdd,ch:models,app:measurement-reference} \\
\addlinespace[0.5em]
Link/system & Does the complete supported combination close? & BER/FEC behavior,
margin waterfall, interoperability, recovery &
\Cref{ch:product-readiness,ch:networking} \\
\addlinespace[0.5em]
Control/management & Can the product enter, report, and recover from required
states? & CMIS state, diagnostics correlation, alarms, restart &
\Cref{sec:bringup,app:measurement-reference} \\
\bottomrule
\end{tabular}
\par\medskip
{\raggedright\footnotesize\textbf{Table~\thetable.} Evidence domains for
Step~5. Instruments and procedures:
\Cref{app:measurement-reference}.\par}
\end{table*}

\textit{Evidence and handoff:} Freeze traceable requirement results, margin
boundaries, and supported combinations; unresolved operating-envelope risk
returns to architecture or characterization before qualification claims rely on
it.

\subsection{Step 6: Qualify reliability}
\label{sec:ladder-stress}

Reliability qualification is a separate evidence discipline within the
product-readiness lifecycle. It is not a more severe version of system
validation. Convert the known operating envelope and credible mechanisms into a
bounded life and environmental confidence argument: mechanism-relevant stresses,
observables, representative samples, acceptance criteria, and confidence.
Distinguish operation while exposed, permanent post-exposure degradation, and
justified life projection. \Cref{ch:reliability} develops the complete method.

\textit{Evidence and handoff:} Accept, derate, or hold life risk for the
envelope; unresolved mechanism ownership returns before unrestricted ramp.

\subsection{Step 7: Validate manufacturing}
\label{sec:ladder-production}

Manufacturing validation does not re-prove system suitability or product
lifetime. Establish that the production reference, measurement system, process
distributions, ATP coverage, genealogy, supplier controls, and reaction plans
can reproduce and control the supported design. A hand-built design may pass
Steps~5 and~6 and still fail here. \Cref{ch:manufacturing} develops the
complete method.

\textit{Evidence and handoff:} Authorize controlled production exposure only
when multi-lot yield, measurement confidence, and escape detection support it.

\subsection{Step 8: Run a controlled pilot}
\label{sec:ladder-fleet}

A pilot is a bounded production-representative deployment with cohort identity,
success criteria, enhanced telemetry, and rollback. It tests whether lab and
factory models survive install practice and traffic mix
(\Cref{sec:evidence-disciplines,ch:networking}).

\textit{Evidence and handoff:} Expand, hold, restrict, or roll back from exit
metrics versus the release model.

\subsection{Step 9: Ramp mass production}
\label{sec:ladder-mp}

Sustain volume under ATP, SPC, supplier, and change controls after pilot exit.
Pilot luck is not proof of sustained control (\Cref{ch:manufacturing,tab:npi}).

\textit{Evidence and handoff:} Increase volume only while factory, supplier, and
early-field indicators remain controlled.

\subsection{Step 10: Monitor the fleet}
\label{sec:ladder-fleet-monitor}

Fleet monitoring compares deployed cohorts with the release model using lane-,
module-, lot-, site-, firmware-, topology-, and installation-age evidence. It
detects distribution shifts and identifies populations requiring containment or
investigation. Fleet correlation prioritizes hypotheses but does not confirm a
mechanism. Average pre-FEC BER can hide whether errors are stationary and sparse
or concentrated into operationally dangerous bursts
(\Cref{sec:fec-symbol-histogram,ch:networking}). Procedures and bucket maps live
in \Cref{sec:fleet-triage,ch:failure-modes}.

\textit{Evidence and handoff:} Contain, sustain, or investigate; return evidence
to the appropriate earlier readiness step when the release model fails.

\subsection{Step 11: Feed learning into the next revision}
\label{sec:ladder-feedback}

Convert fleet and manufacturing evidence into changed requirements, designs,
qualification, ATP, or controls (\Cref{sec:systems-engineering-loop}).

\textit{Evidence and handoff:} Write next-revision targets with owners, or
explicitly accept residual risk without change.

\section{Choosing evidence within a step}
\label{sec:ladder-evidence-class}

An experiment is not inherently characterization, verification, validation, or
qualification. Its category is determined by the claim being made
(\Cref{sec:evidence-disciplines}). Use the lowest-cost evidence that changes the
decision. State reference plane, condition, population, uncertainty, and
applicability. Reuse evidence across claims only when those conditions support
the new claim.

A first discriminating measurement should separate broad ownership classes; for
an optical-link symptom, named-plane power is often an inexpensive first split
between gross attenuation and signal-quality hypotheses
(\Cref{sec:validation-fork,sec:debug-fork}).

\heuristic{Name the reference plane before you name the instrument. A pretty
eye at the wrong plane is a wrong answer.}

\keyidea{Product readiness is a sequence of decisions, not a catalog of tests.
Each step should produce evidence for one explicit decision and identify the
uncertainty that remains.}

\section{Bring-up and production-representative corners}
\label{sec:bringup}

Bring-up proves a module and system can be powered, managed, and linked the way
production will run them. A quiet room-temperature BERT result is not system
validation.

\begin{enumerate}
\item Confirm identity, rails, temperature, and firmware.
\item Reach the required management-ready state.
\item Enable the intended lanes or light under host control.
\item Establish the optical path and verify basic power.
\item Establish electrical lock and lane alignment.
\item Run basic traffic and capture a reproducible baseline.
\end{enumerate}

A module forced into an emitting state and passing BER has not passed bring-up
if its required management sequence, safe state, diagnostics, alarms, and
recovery behavior are incorrect. Register-level CMIS detail lives in
\Cref{sec:cmis,app:measurement-reference}.

\paragraph{Production-representative corners.}
\label{sec:prod-corners}

\begin{table*}[ht]
\refstepcounter{table}\label{tab:prod-corners}%
\centering
\setlength{\tabcolsep}{3pt}%
\footnotesize
\begin{tabular}{@{}p{0.28\linewidth}p{0.66\linewidth}@{}}
\toprule
Corner class & Representative challenges \\
\midrule
Thermal and loading & Case temperature, airflow, neighboring modules, full
traffic \\
\addlinespace[0.5em]
Electrical and host & Supported hosts, voltage corners, SerDes/equalization,
reset and restart \\
\addlinespace[0.5em]
Optical plant & Production fiber, connectors, reflections, loss, supported peers \\
\addlinespace[0.5em]
Control and service & CMIS transitions, alarms, hot-swap, recovery, firmware
combinations \\
\bottomrule
\end{tabular}
\par\medskip
{\raggedright\footnotesize\textbf{Table~\thetable.} Production-representative
corners for system validation. Mechanism detail:
\Cref{ch:lasers,ch:wdm,ch:networking}.\par}
\end{table*}

\label{sec:prod-corners-read}%

\section{Interview takeaway}
\label{sec:validation-design-review}

\keyidea{Staff-level readiness leadership is not the ability to name many tests.
It is the ability to state which uncertainty remains, which evidence removes it,
which decision follows, and which residual risk moves into the next lifecycle
step. Verification, system validation, reliability qualification, manufacturing
validation, and fleet monitoring are connected evidence streams, not
interchangeable labels.}

Junior mistake: using EVT, DVT, PVT, verification, validation, qualification,
and production test as interchangeable names for ``testing.'' Better practice:
name the claim, population, condition, evidence, and decision. Then use the
engineering term that matches that claim
(\Cref{tab:ladder,tab:validation-jobs,ch:manufacturing,app:decision-trees}).

\subsection{Interview Q\&A: Optical Product Readiness}
\label{sec:validation-interview-qa}

Practice speaking these answers aloud. Prefer first-person reasoning over
definitions. Detail lives earlier in this chapter
(\Cref{sec:product-readiness,tab:ladder,sec:evidence-disciplines}).
Score your answer using the chapter-end spoken-answer rubric
(\Cref{sec:chapter-interview-rubric}).

\paragraph{Question 1. What is product readiness, and where does system
validation fit?}
\textit{Tests:} umbrella versus discipline.

\textit{Spoken answer.}
``I use product readiness for the complete lifecycle from requirements through
fleet learning. Within that lifecycle, system validation demonstrates intended
use, reliability qualification addresses permanent time- and exposure-driven
degradation, and manufacturing validation establishes factory reproducibility.''

\textit{Pressure follow-up.} ``When is system validation sufficient?''\\
\textit{Answer pivot.} ``When the supported operating and interoperability
envelope is closed for the claim under test. That is not unrestricted production
readiness.''

\textit{Trap:} treating validation as the name of the entire readiness program.

\paragraph{Question 2. What is the difference between characterization,
verification, system validation, reliability qualification, manufacturing
validation, production acceptance testing, and fleet monitoring?}
\textit{Tests:} terminology discipline.

\textit{Spoken answer.}
``Characterization maps behavior. Verification checks a frozen requirement at a
named plane. System validation asks whether the complete product works for
intended use. Reliability qualification bounds permanent degradation.
Manufacturing validation asks whether production can reproduce and control the
design. ATP dispositions a unit or population. Fleet monitoring asks whether the
deployed population matches the release model''
(\Cref{tab:validation-jobs}).

\textit{Pressure follow-up.} ``Where does burn-in fit?''\\
\textit{Answer pivot.} ``Burn-in is a production screen when justified. It does
not replace life qualification.''

\textit{Trap:} calling them ``levels of testing.''

\paragraph{Question 3. Walk me through the optical product-readiness lifecycle.
Why is it ordered, and which activities can overlap?}
\textit{Tests:} eleven steps, order, and overlap.

\textit{Spoken answer.}
``Define requirements, review architecture, bring up hardware, characterize
behavior, verify requirements and validate system use, qualify reliability,
validate manufacturing, run a controlled pilot, ramp mass production, monitor
the fleet, and feed learning forward. Calendar work can overlap, but a later
pass cannot substitute for missing earlier evidence'' (\Cref{tab:ladder}).

\textit{Pressure follow-up.} ``Where do EVT, DVT, and PVT fit?''\\
\textit{Answer pivot.} ``They are program-phase labels, not evidence
definitions. Translate each gate into hardware population, required evidence,
open risks, and exit decision.''

\textit{Trap:} treating EVT/DVT/PVT as universal engineering definitions.

\paragraph{Question 4. What happens during architecture review?}
\textit{Tests:} budget and derating priors.

\textit{Spoken answer.}
``I ask whether the design can plausibly meet requirements before hardware makes
changes expensive. I close the main budgets on stated assumptions and name the
thin margins later evidence must challenge.''

\textit{Pressure follow-up.} ``Where do derating rules come from?''\\
\textit{Answer pivot.} ``Priors from vendor data, prior measurements,
physics-of-failure models, and field history. Qualification later checks this
design'' (\Cref{ch:reliability}).

\textit{Trap:} ``I simulated the link budget and it closed.''

\paragraph{Question 5. What is the objective of hardware bring-up?}
\textit{Tests:} reproducible baseline.

\textit{Spoken answer.}
``Bring-up establishes a known, reproducible operating state: identity, rails,
firmware, management-ready state, light, lock, and a simple link. I separate
integration fails from product-performance questions'' (\Cref{sec:bringup}).

\textit{Pressure follow-up.} ``What if ready state never appears?''\\
\textit{Answer pivot.} ``Freeze the setup, dump management state and rails, and
isolate host versus module before deep optical sweeps.''

\textit{Trap:} ``power on and check BER.''

\paragraph{Question 6. What is characterization trying to produce?}
\textit{Tests:} behavioral model.

\textit{Spoken answer.}
``A behavioral model: nominals, distributions, sensitivities, and signatures.
The output is which margins are thin and which corners to challenge in system
validation, not only pass or fail.''

\textit{Pressure follow-up.} ``Every unit versus sample?''\\
\textit{Answer pivot.} ``Cheap identity and power checks can be every-unit;
expensive waterfalls stay sample or audit unless escape data forces otherwise.''

\textit{Trap:} ``measure over temperature.''

\paragraph{Question 7. Why are production-representative corners required?}
\textit{Tests:} quiet-bench versus shipping envelope.

\textit{Spoken answer.}
``A room-temperature golden-host BERT result is not system validation. I need
thermal and loading, electrical and host, optical plant, and control and service
corners that match the supported envelope'' (\Cref{tab:prod-corners}).

\textit{Pressure follow-up.} ``What is the cheapest corner that often moves
BER?''\\
\textit{Answer pivot.} ``Loaded chassis thermal with production fiber and ORL,
because faceplate temperature and reflections rarely match a quiet bench.''

\textit{Trap:} treating chamber setpoints as case temperature.

\paragraph{Question 8. In Step~5, what is the difference between verifying
requirements and margin versus validating interoperability and intended use?}
\textit{Tests:} verification versus system validation.

\textit{Spoken answer.}
``Requirement and margin verification closes frozen specs at named planes.
Intended-use validation asks whether production-representative hosts, peers,
firmware, and fiber plants keep those boundaries intact''
(\Cref{sec:ladder-verify,sec:ladder-system-validation}).

\textit{Pressure follow-up.} ``Why not test only the reference host?''\\
\textit{Answer pivot.} ``Reference-host margin can look fine while a production
host EQ or CMIS path moves the cliff.''

\textit{Trap:} treating a verified requirement as complete system validation.

\paragraph{Question 9. What must accompany every readiness measurement?}
\textit{Tests:} plane and conditions.

\textit{Spoken answer.}
``Reference plane, metric definition, population, condition, uncertainty, and
the decision the result enables. A number without those is not evidence''
(\Cref{sec:ladder-evidence-class,app:measurement-reference}).

\textit{Pressure follow-up.} ``Supplier quotes $-10$~dBm sensitivity. First
question?''\\
\textit{Answer pivot.} ``Average power or OMA, at which plane and BER or FEC
condition?''

\textit{Trap:} comparing numbers from different planes.

\paragraph{Question 10. How do you choose the next readiness measurement?}
\textit{Tests:} information per cost.

\textit{Spoken answer.}
``I pick the cheapest measurement that separates the strongest surviving
hypotheses and changes the decision. Named-plane power often splits gross
attenuation from signal-quality hypotheses before deeper instruments.''

\textit{Pressure follow-up.} ``When is destructive analysis justified?''\\
\textit{Answer pivot.} ``When non-destructive evidence cannot separate owners and
the release or containment decision is blocked.''

\textit{Trap:} ``run full optical characterization.''

\paragraph{Question 11. Why is pilot deployment part of the product-readiness
lifecycle?}
\textit{Tests:} bounded field experiment.

\textit{Spoken answer.}
``A pilot is Step~8 operational readiness evidence. Bounded serials, enhanced
telemetry, exit criteria, and rollback. Compare BER, FEC behavior, retrains,
temperature, power, and cohort rates to the release model''
(\Cref{sec:ladder-fleet,ch:networking}).

\textit{Pressure follow-up.} ``What justifies expanding the pilot?''\\
\textit{Answer pivot.} ``Exit metrics met, no unexplained cohort, and
containment still reversible.''

\textit{Trap:} calling a small shipment a pilot.

\paragraph{Question 12. Give me a 60-second answer for establishing readiness
for a new optical module.}
\textit{Tests:} time-boxed program.

\textit{Spoken answer.}
``Define measurable requirements and the release call. Review whether
architecture budgets close. Bring up reproducibly, characterize distributions
and cliffs, then verify requirements and validate intended use on
production-like hosts and plants. Qualify named life mechanisms and validate
manufacturing measurement, ATP, and yield. Run a controlled pilot, ramp under
monitoring, and feed escapes back into requirements or controls.''

\textit{Pressure follow-up.} ``Schedule is cut in half. What do you protect?''\\
\textit{Answer pivot.} ``Requirements, bring-up integrity, the thinnest margin
corners, and a reversible pilot.''

\textit{Trap:} listing BER, temperature, reliability, and interop with no order.
